\section{Segment Tree iterativo}
\whenuse{Arreglo din\'amico con updates puntuales y consultas de una operaci\'on
asociativa sobre rangos; tambi\'en para buscar la primera posici\'on que cumple.}

Versi\'on para suma, actualizaci\'on puntual y consultas semiabiertas $[l,r)$.
Cambiar \texttt{combine} y \texttt{NEUTRAL} para min, max, gcd, etc.

\textbf{Operaciones comunes}
\begin{center}
\begin{tabular}{|c|c|c|}
\hline
Consulta & \texttt{oper} & \texttt{NEUT} \\
\hline
Minimo & \texttt{min} & \texttt{INF} \\
Maximo & \texttt{max} & \texttt{-INF} \\
Suma & \texttt{+} & \texttt{0} \\
Producto & \texttt{*} & \texttt{1} \\
XOR & \texttt{\^} & \texttt{0} \\
OR bit a bit & \texttt{|} & \texttt{0} \\
AND bit a bit & \texttt{\&} & \texttt{(\textasciitilde 0)} \\
GCD & \texttt{gcd} & \texttt{0} \\
LCM & \texttt{lcm} & \texttt{1} \\
\hline
\end{tabular}
\end{center}

\complexity{build $O(n)$; update/query $O(\log n)$; memoria $O(n)$.}

\begin{lstlisting}[style=cpp]
struct SegmentTree {
    int n = 1;
    vector<ll> tree;

    explicit SegmentTree(int size) {
        while (n < size) n <<= 1;
        tree.assign(2 * n, 0);
    }

    explicit SegmentTree(const vector<ll> &a) : SegmentTree(a.size()) {
        copy(a.begin(), a.end(), tree.begin() + n);
        for (int p = n - 1; p; --p)
            tree[p] = tree[p << 1] + tree[p << 1 | 1];
    }

    void set_value(int pos, ll value) {
        int p = pos + n;
        tree[p] = value;
        while (p >>= 1)
            tree[p] = tree[p << 1] + tree[p << 1 | 1];
    }

    ll query(int l, int r) const { // [l, r)
        ll left = 0, right = 0;
        for (l += n, r += n; l < r; l >>= 1, r >>= 1) {
            if (l & 1) left += tree[l++];
            if (r & 1) right = tree[--r] + right;
        }
        return left + right;
    }
};
\end{lstlisting}
